%
%      Brno University of Technology
%      Faculty of Information Technology
%
%      MSc Thesis	2009/2010
%
%      Design of S-Boxes Using Genetic Algorithms
%
%
%      Bedrich Hovorka
%
Tato příloha vysvětluje značení v chromozomu kartézského genetického programování.
Notace byla použita ve shodě s nástrojem CGPTools \cite{vasicek}.
\subsubsection{Tvar chromozomu}
Hlavička obsahuje údaje v tomto pořadí:
\begin{verbatim}
  {počet vstupů, počet výstupů, sloupce, řádky, počet vstupů do bloku, l-back,
   počet použitých bloků}
\end{verbatim}
Následují popisy jednotlivých bloků (za vstupní blok se považuje i vstup celého obvodu; bloky
se číslují od počtu vstupů):
\begin{verbatim}
  ([číslo bloku] vstupní blok a, vstupní blok b, funkce)
\end{verbatim}
Na závěr jsou výstupy některých bloků (uvedená čísla) prohlášeny za výstupy celého obvodu:
\begin{verbatim}
  (seznam bloků)
\end{verbatim}


\subsubsection{Přiřazení funkcí}
Pro vstupní proměnné $a$ a $b$:

\begin{center}
     \begin{tabular}{| c | c |}
      \hline
      Číslo funkce & Výraz \\ \hline\hline
      0 & $a$ \\ \hline
      1 & $a$ AND $b$ \\ \hline
      2 & $a$ OR $b$ \\ \hline
      3 & $a$ XOR $b$ \\ \hline
      4 & NOT $a$ \\ \hline
      5 & NOT $b$ \\ \hline
      6 & $a$ AND (NOT $b$) \\ \hline
      7 & $a$ NAND $b$ \\ \hline
      8 & $a$ NOR $b$ \\ \hline
     \end{tabular}
\end{center}
